\documentclass[11pt]{article}
\usepackage{notes}
\definecolor{AccN}{HTML}{6D4C41}
\begin{document}
\titlepg{6}{Neural Networks:\\Hidden Layers and Backpropagation}{Part I: why hidden layers --- the single neuron, the XOR wall, non-linear activations, learned feature maps and universal approximation. Part II: backpropagation by hand --- one forward and one backward pass through a 2--2--1 sigmoid network}
\tableofcontents
\newpage
\section{How to Use This Text}
This reference text accompanies the Lecture~6 slides. It has two parts that answer the two questions every analyst should be able to answer about a neural network.
\begin{itemize}
\item \textbf{Part I --- Why hidden layers?} What can a single artificial neuron do, what can it never do, and why does a \emph{hidden} layer with a \emph{non-linear} activation remove that limit? The answer turns on one small problem, XOR, and one big idea: hidden layers are \emph{learned feature maps}.
\item \textbf{Part II --- Backpropagation by hand.} Given that hidden layers are needed, how are their weights learned? One example is pushed forward and backward through a 2--2--1 network, with every number computed.
\end{itemize}
Coloured boxes mark the different kinds of material: \textbf{Definition} boxes give terms you must be able to state; \textbf{Derivation} boxes show where formulas come from; \textbf{Worked example} boxes carry complete calculations; \textbf{Defence \& wargaming} boxes connect ideas to practice; \textbf{Pitfall} boxes describe common errors; \textbf{Check your understanding} boxes pose questions whose answers follow.

Every network in the rest of the course is built from the ideas of this lecture: the convolutional network of Lecture~7, the transformers of Lectures~8 and~9, and the deep Q-networks and policy networks of Lectures~10--12. Lecture~1 introduced neurons and backpropagation at headline level; this lecture does the mathematics.

\part{Why Hidden Layers?}

\section{The Artificial Neuron}
\subsection{Definition}
\begin{definition}{artificial neuron}
An artificial neuron with inputs $x=(x_1,\dots,x_d)$ computes a \textbf{pre-activation} $z=w\cdot x+b=\sum_{i=1}^dw_ix_i+b$ and an \textbf{output} $a=f(z)$. The numbers $w_i$ are \textbf{weights}, $b$ is the \textbf{bias}, and $f$ is the \textbf{activation function}.
\end{definition}
The weights say how much each input matters and in which direction; the bias shifts the point at which the neuron ``switches on''. Logistic regression (Lecture~2) is exactly one neuron with a sigmoid activation, $f(z)=1/(1+e^{-z})$, read as a probability.

\subsection{One neuron draws one straight line}
To classify, we threshold the output: predict class 1 when $f(z)$ exceeds some level $\tau$. For every common activation --- the step, the sigmoid, tanh, ReLU on its active side --- $f$ is \emph{monotonic increasing}, so $f(z)>\tau$ is equivalent to $z>c$ for some constant $c$ (for the sigmoid with $\tau=0.5$, $c=0$). Absorbing $c$ into the bias, the decision rule is always
\[\hat y=1\iff w_1x_1+w_2x_2+b>0.\]
\begin{derivation}{the decision boundary is a straight line}
The boundary is the set where $w_1x_1+w_2x_2+b=0$. If $w_2\neq0$ this rearranges to
\[x_2=-\frac{w_1}{w_2}\,x_1-\frac{b}{w_2},\]
a straight line with slope $-w_1/w_2$ and intercept $-b/w_2$. The weight vector $w=(w_1,w_2)$ is perpendicular to the line and points towards the class-1 side. In $d$ dimensions the boundary $w\cdot x+b=0$ is a flat \emph{hyperplane}.
\end{derivation}
Training a single neuron can move and rotate this line, but never bend it. Problems whose classes can be split by one straight line (one hyperplane) are called \textbf{linearly separable}.

\begin{example}{AND and OR with one neuron}
Use a step output, $\hat y=1$ if $z>0$.
\begin{center}\small\renewcommand{\arraystretch}{1.2}
\begin{tabular}{cc|rc|rc}\toprule
& & \multicolumn{2}{c|}{AND: $w=(1,1)$, $b=-1.5$} & \multicolumn{2}{c}{OR: $w=(1,1)$, $b=-0.5$}\\
$x_1$ & $x_2$ & $z$ & $\hat y$ & $z$ & $\hat y$\\\midrule
0&0&$-1.5$&0&$-0.5$&0\\
0&1&$-0.5$&0&$0.5$&1\\
1&0&$-0.5$&0&$0.5$&1\\
1&1&$0.5$&1&$1.5$&1\\\bottomrule
\end{tabular}\end{center}
The AND boundary is the line $x_1+x_2=1.5$; the OR boundary is $x_1+x_2=0.5$. Both are correct on all four inputs. Changing only the bias moves the line parallel to itself and turns AND into OR.
\end{example}

\subsection{The perceptron learning rule}
Rosenblatt's perceptron (1958) learned such lines from examples. For each training example with label $y\in\{0,1\}$ and prediction $\hat y$, it updates
\[w\leftarrow w+\eta\,(y-\hat y)\,x,\qquad b\leftarrow b+\eta\,(y-\hat y).\]
Nothing changes when the prediction is right; a missed positive pulls the line towards the example, a false alarm pushes it away. The \textbf{perceptron convergence theorem} guarantees that this finds a separating line in finitely many steps \emph{if one exists}.
\begin{example}{learning AND with the perceptron rule}
Start from $w=(0,0)$, $b=0$, $\eta=1$, step output with $\hat y=1$ iff $z>0$. Visit the examples in the order $(0,0),(0,1),(1,0),(1,1)$.\\
Epoch 1: $(0,0)$: $z=0\Rightarrow\hat y=0$, correct. $(0,1)$: $z=0$, correct. $(1,0)$: correct. $(1,1)$: $z=0\Rightarrow\hat y=0$, but $y=1$, so $w=(1,1)$, $b=1$.\\
Epoch 2: $(0,0)$: $z=1\Rightarrow\hat y=1$, wrong: $b=0$. $(0,1)$: $z=1$, wrong: $w=(1,0)$, $b=-1$. $(1,0)$: $z=0$, correct. $(1,1)$: $z=0$, wrong: $w=(2,1)$, $b=0$.\\
Epoch 3: $(0,0)$: $z=0$, correct. $(0,1)$: $z=1$, wrong: $w=(2,0)$, $b=-1$. $(1,0)$: $z=1$, wrong: $w=(1,0)$, $b=-2$. $(1,1)$: $z=-1$, wrong: $w=(2,1)$, $b=-1$.\\
Epoch 4: $(0,0)$: $z=-1$, correct. $(0,1)$: $z=0$, correct. $(1,0)$: $z=1$, wrong: $w=(1,1)$, $b=-2$. $(1,1)$: $z=0$, wrong: $w=(2,2)$, $b=-1$.\\
Epoch 5: $(0,0)$: $z=-1$, correct. $(0,1)$: $z=1$, wrong: $w=(2,1)$, $b=-2$. $(1,0)$: $z=0$, correct. $(1,1)$: $z=1$, correct.\\
Epoch 6: $z=-2,\,-1,\,0,\,1$ for the four inputs: all correct, and training stops with $w=(2,1)$, $b=-2$, the line $2x_1+x_2=2$. It differs from the hand-chosen $x_1+x_2=1.5$ but separates the data equally well: a separable problem has many solutions.
\end{example}

\section{The XOR Wall}\label{sec:xor}
\subsection{The problem}
\textbf{Exclusive OR} (XOR) outputs 1 when exactly one input is 1. Plotted on the unit square, class 1 sits at $(0,1)$ and $(1,0)$ and class 0 at $(0,0)$ and $(1,1)$: the classes occupy opposite corners. Try any straight line: whichever side contains both class-1 corners also contains at least one class-0 corner.

\begin{center}\begin{tikzpicture}[scale=2.2]
\draw[GreyN!40](0,0)grid[step=0.5](1,1);
\draw[arr](-0.1,0)--(1.3,0)node[right]{$x_1$};\draw[arr](0,-0.1)--(0,1.3)node[above]{$x_2$};
\foreach \p in {(0,0),(1,1)}{\node[circle,draw=BluN,fill=BluN!30,inner sep=2.5pt]at\p{};}
\foreach \p in {(0,1),(1,0)}{\node[rectangle,draw=RedN,fill=RedN!35,inner sep=3pt]at\p{};}
\draw[RedN,dashed,thick](-0.1,0.55)--(1.2,0.7);\node[font=\small,RedN]at(0.8,0.9){any line fails};
\end{tikzpicture}\\*{\small\color{GreyN}Figure 1: XOR. Squares are class 1, circles class 0.}\end{center}

\subsection{Proof that no single neuron computes XOR}
\begin{derivation}{XOR is not linearly separable}
Suppose weights $w_1,w_2$ and bias $b$ existed such that the output is 1 exactly when $z=w_1x_1+w_2x_2+b>0$. The four rows of the truth table require
\[\begin{array}{lll}
(0,0)\mapsto0: & b\le0 & (1)\\
(1,0)\mapsto1: & w_1+b>0 & (2)\\
(0,1)\mapsto1: & w_2+b>0 & (3)\\
(1,1)\mapsto0: & w_1+w_2+b\le0 & (4)
\end{array}\]
Adding (2) and (3) gives $w_1+w_2+2b>0$, so $w_1+w_2+b>-b$. By (1), $-b\ge0$, so $w_1+w_2+b>0$, contradicting (4). No such weights exist. \hfill$\blacksquare$
\end{derivation}
The proof does not depend on which monotonic activation is used, so the same impossibility holds for logistic regression, a sigmoid neuron or a ReLU neuron.

\subsection{What logistic regression actually does on XOR}\label{sec:logxor}
It is instructive to see what a single sigmoid neuron \emph{learns} when trained on XOR with the cross-entropy loss.
\begin{example}{the best single neuron is useless}
Take $w=(0,0)$, $b=0$: every input gets $p=\sigma(0)=0.5$. The cross-entropy gradient of a sigmoid neuron is $\sum_n(p_n-y_n)x_n$ for the weights and $\sum_n(p_n-y_n)$ for the bias (derived in Part~II). Here
\begin{align*}\textstyle\sum_n(p_n-y_n)x_n&=(0.5)(0,0)+(-0.5)(0,1)+(-0.5)(1,0)+(0.5)(1,1)=(0,0),\\ \textstyle\sum_n(p_n-y_n)&=0.5-0.5-0.5+0.5=0.\end{align*}
The gradient is zero. Because the logistic-regression loss is convex, this stationary point is the \emph{global} minimum: the best single neuron predicts 0.5 for every input, with loss $\ln2\approx0.693$ per example. It has learned nothing.
\end{example}

\subsection{History}
Minsky and Papert's book \emph{Perceptrons} (1969) made this limitation famous. Multi-layer networks could in principle solve XOR, but nobody had a practical way to train their hidden weights, and research funding moved elsewhere. The revival came when Rumelhart, Hinton and Williams (1986) popularised \textbf{backpropagation} --- the subject of Part~II.

\begin{defence}{an XOR-shaped threat picture}
Consider air tracks described by speed and altitude (an illustrative, simplified picture). Fast-and-low tracks match sea-skimming or cruise-missile profiles (hostile); slow-and-high tracks match loitering surveillance platforms (hostile); slow-and-low tracks are often own helicopters, and fast-and-high tracks airliners or own jets (routine). Neither speed alone nor altitude alone separates hostile from routine; only their \emph{combination} does --- the XOR pattern. A single-neuron classifier on these two inputs can do no better than the majority class. Real data contain many such interactions: ``either A or B but not both'', a target inside a ring of clutter, interleaved patterns.
\end{defence}

\section{Width and Linear Depth Are Not Enough}
\subsection{One wide layer is still linear}
A layer of $K$ neurons computes $\hat y_k=f(w_k\cdot x+b_k)$ for $k=1,\dots,K$. Each output draws its own straight line on the \emph{same raw inputs}, and no output sees another's answer. For XOR every one of the $K$ neurons faces the same impossibility proof. With a softmax over the outputs (multinomial logistic regression), class $k$ is predicted where $w_k\cdot x+b_k$ is largest; the regions are separated by straight lines and each region is \emph{convex} --- so no class can occupy two opposite corners. Width adds lines but cannot \emph{reason with} them (``inside line 1 AND outside line 2''). That requires a second stage whose inputs are the first stage's outputs.

\subsection{Stacked linear layers collapse into one}\label{sec:collapse}
\begin{derivation}{composition of linear maps}
Two layers with the identity activation compute
\[\hat y=W_2(W_1x+b_1)+b_2=\underbrace{(W_2W_1)}_{W}x+\underbrace{(W_2b_1+b_2)}_{b}.\]
The product of two matrices is a matrix and the rest is a vector, so the two-layer network is exactly a one-layer network with weights $W$ and bias $b$. By induction, any number of linear layers collapses into one.
\end{derivation}
\begin{example}{checking the collapse numerically}
$W_1=\begin{pmatrix}1&2\\0&1\end{pmatrix}$, $b_1=\begin{pmatrix}1\\0\end{pmatrix}$, $W_2=(1\ \ -1)$, $b_2=0.5$. Collapsed: $W=W_2W_1=(1\cdot1+(-1)\cdot0,\ 1\cdot2+(-1)\cdot1)=(1,\ 1)$ and $b=W_2b_1+b_2=1+0+0.5=1.5$.\\
At $x=(1,1)$: the two-layer network gives $h=(1+2+1,\ 1)=(4,1)$ and $\hat y=4-1+0.5=3.5$; the collapsed layer gives $1+1+1.5=3.5$.\\
At $x=(2,-1)$: $h=(2-2+1,\ -1)=(1,-1)$, $\hat y=1+1+0.5=2.5$; collapsed: $2-1+1.5=2.5$.
\end{example}
\begin{keypoint}
Depth without non-linearity is an illusion. The hidden layer must \emph{bend} its output, or the whole network remains one straight boundary. This is why every convolution in a CNN (Lecture~7) and every feed-forward block in a transformer (Lecture~9) is followed by a non-linear activation.
\end{keypoint}

\section{Non-Linear Activation Functions}
\begin{center}\small\renewcommand{\arraystretch}{1.35}
\begin{tabular}{llll}\toprule
\textbf{Name} & $f(z)$ & $f'(z)$ & \textbf{Notes}\\\midrule
Step & $1$ if $z>0$, else $0$ & $0$ (undefined at 0) & perceptron; no usable gradient\\
Sigmoid & $1/(1+e^{-z})$ & $\sigma(z)(1-\sigma(z))\le0.25$ & range $(0,1)$; probability output\\
Tanh & $\tanh z$ & $1-\tanh^2z\le1$ & range $(-1,1)$; zero-centred\\
ReLU & $\max(0,z)$ & $1$ if $z>0$, $0$ if $z<0$ & default for hidden layers\\\bottomrule
\end{tabular}\end{center}
\begin{center}\begin{tikzpicture}
\begin{axis}[width=10cm,height=5.4cm,xmin=-3,xmax=3,ymin=-1.2,ymax=2.2,grid=major,grid style={GreyN!20},xlabel={$z$},legend pos=outer north east,legend style={font=\footnotesize}]
\addplot[GreyN,line width=1pt,domain=-3:0,samples=2]{0};\addlegendentry{step}
\addplot[BluN,line width=1.1pt,domain=-3:3,samples=60]{1/(1+exp(-x))};\addlegendentry{sigmoid}
\addplot[GrnN,line width=1.1pt,domain=-3:3,samples=60]{tanh(x)};\addlegendentry{tanh}
\addplot[RedN,line width=1.1pt,domain=-3:3,samples=61]{max(0,x)};\addlegendentry{ReLU}
\addplot[GreyN,line width=1pt,domain=0:3,samples=2,forget plot]{1};
\end{axis}\end{tikzpicture}\\*{\small\color{GreyN}Figure 2: four activation functions.}\end{center}
The only requirement for breaking the collapse is that $f$ is \emph{not} of the form $az+b$. ReLU --- two straight pieces joined at a kink --- is enough. It is preferred in hidden layers because it is cheap and its derivative is exactly 1 for active units, so gradients do not shrink as they pass back through many layers; the sigmoid's derivative is at most 0.25, which causes the \emph{vanishing-gradient} problem examined in Part~II. The step function has zero derivative almost everywhere, so gradient-based training cannot use it.

\section{A Hidden Layer Solves XOR}
\subsection{The network}
Take two hidden ReLU units and a linear output:
\[h=\text{ReLU}(Wx+c),\qquad\hat y=w^{\mathsf T}h,\qquad W=\begin{pmatrix}1&1\\1&1\end{pmatrix},\ c=\begin{pmatrix}0\\-1\end{pmatrix},\ w=\begin{pmatrix}1\\-2\end{pmatrix}.\]
So $h_1=\text{ReLU}(x_1+x_2)$, $h_2=\text{ReLU}(x_1+x_2-1)$ and $\hat y=h_1-2h_2$. This is the classic solution in Goodfellow, Bengio and Courville, \emph{Deep Learning} (2016), \S6.1.
\begin{example}{the forward pass for all four inputs}
\begin{center}\small\renewcommand{\arraystretch}{1.25}
\begin{tabular}{cc|cc|cc|c|c}\toprule
$x_1$&$x_2$&$z_1$&$z_2$&$h_1$&$h_2$&$\hat y=h_1-2h_2$&XOR\\\midrule
0&0&0&$-1$&0&0&0&0\\
0&1&1&0&1&0&1&1\\
1&0&1&0&1&0&1&1\\
1&1&2&1&2&1&0&0\\\bottomrule
\end{tabular}\end{center}
All four outputs are correct. In row 1, ReLU clips $z_2=-1$ to 0; without the clip, $h_2=-1$ and $\hat y=0-2(-1)=2$, which is wrong.
\end{example}
\subsection{Reading the hidden units}
$h_1$ counts \emph{how many} inputs are on (an ``OR-like'' detector); $h_2$ is non-zero only when \emph{both} are on (an ``AND'' detector). The output says ``at least one, but subtract heavily if both'': XOR $=$ OR AND NOT AND. The hidden layer computes the pieces; the output combines them --- the combining that a single layer could not do.
\begin{example}{the same network without ReLU}
Replace ReLU by the identity: $\hat y=z_1-2z_2=(x_1+x_2)-2(x_1+x_2-1)=-x_1-x_2+2$. The outputs for $(0,0),(0,1),(1,0),(1,1)$ are $2,1,1,0$. The network has collapsed to one straight line (one linear function), exactly as Section~\ref{sec:collapse} predicts, and it no longer computes XOR.
\end{example}

\section{Hidden Layers as Learned Feature Maps}
\subsection{Re-mapping the data}
The hidden layer is a function $x\mapsto h(x)$ that sends each input point to a new point. For the XOR network:
\[(0,0)\mapsto(0,0),\qquad(0,1)\mapsto(1,0),\qquad(1,0)\mapsto(1,0),\qquad(1,1)\mapsto(2,1).\]
\begin{center}\begin{tikzpicture}[scale=1.6]
\begin{scope}
\draw[GreyN!40](0,0)grid[step=0.5](1,1);\draw[arr](-0.1,0)--(1.3,0)node[right]{$x_1$};\draw[arr](0,-0.1)--(0,1.3)node[above]{$x_2$};
\foreach \p in {(0,0),(1,1)}{\node[circle,draw=BluN,fill=BluN!30,inner sep=2.5pt]at\p{};}
\foreach \p in {(0,1),(1,0)}{\node[rectangle,draw=RedN,fill=RedN!35,inner sep=3pt]at\p{};}
\node[font=\small]at(0.5,-0.4){input space: not separable};
\end{scope}
\draw[arr,AccN,line width=1.2pt](1.7,0.5)--node[above]{$h$}(2.5,0.5);
\begin{scope}[xshift=3cm]
\draw[GreyN!40](0,0)grid[step=0.5](2,1);\draw[arr](-0.1,0)--(2.3,0)node[right]{$h_1$};\draw[arr](0,-0.1)--(0,1.3)node[above]{$h_2$};
\node[circle,draw=BluN,fill=BluN!30,inner sep=2.5pt]at(0,0){};\node[circle,draw=BluN,fill=BluN!30,inner sep=2.5pt]at(2,1){};
\node[rectangle,draw=RedN,fill=RedN!35,inner sep=3pt]at(1,0){};\node[font=\scriptsize,RedN,above right]at(1,0){$\times2$};
\draw[GrnN,line width=1.2pt](0.4,-0.15)--(2.15,0.72);\node[font=\small]at(1,-0.4){hidden space: one line separates};
\end{scope}
\end{tikzpicture}\\*{\small\color{GreyN}Figure 3: the hidden layer moves both class-1 points to the same spot.}\end{center}
In the hidden coordinates both class-1 inputs land on the same point $(1,0)$, and a single straight line --- the output neuron, $h_1-2h_2=0.5$ --- separates the classes. The output neuron is still only a line, but a line \emph{in the new space}.
\begin{definition}{feature map and representation learning}
The activations of a hidden layer, viewed as a function of the input, form a \textbf{feature map}: a learned change of coordinates. \textbf{Representation learning} means choosing this map by training, so that the final problem becomes easy (often linearly separable), instead of having an engineer hand-craft the features.
\end{definition}
In a dense layer, a hidden unit's feature is one number per example. In a CNN (Lecture~7) each filter's feature map is a whole image --- the filter's response at every location. In a transformer (Lectures~8--9) the features are context-aware vectors for each token.

\subsection{Advantages of learned feature maps}
\begin{itemize}
\item \textbf{Automatic feature engineering:} the network discovers useful intermediate quantities (``how many on?'', ``both on?'') from labels alone.
\item \textbf{Re-usable:} the same features serve many outputs and tasks; this is the basis of transfer learning.
\item \textbf{Compact:} a few hundred features can summarise thousands of raw inputs.
\item \textbf{Inspectable:} feature maps can be visualised to see what the network responds to (with care --- see below).
\item \textbf{Hierarchical:} deeper layers combine shallower features into parts and objects.
\end{itemize}

\begin{defence}{hidden units as learned indicators}
On the air-track data of Section~\ref{sec:xor}, a trained hidden layer could develop units that fire for regions of the speed--altitude plane: ``fast and low'', ``slow and high'', ``slow and low'', ``fast and high''. The output then needs only a weighted vote --- positive weights on the two hostile indicators, negative on the two routine ones --- which is linear in the indicator space. Add radar cross-section, heading-change rate and IFF response as inputs, and the hidden layer can learn combinations no analyst has written down. \textbf{Reality check:} real hidden units rarely align with neat human labels; features are usually spread across many units. Interpreting them needs care, and high-consequence calls need a human analyst.
\end{defence}

\section{How Far Can Hidden Layers Go?}
\subsection{Each ReLU adds a bend}
A ReLU network computes a continuous \emph{piecewise-linear} function: each hidden unit contributes one kink. Combining three units gives a bump:
\begin{example}{a bump from three ReLUs}
$f(x)=\text{ReLU}(x)-2\,\text{ReLU}(x-1)+\text{ReLU}(x-2)$.\\
For $x\le0$ all terms vanish: $f=0$. For $0\le x\le1$: $f=x$, so $f(0.5)=0.5$, $f(1)=1$. For $1\le x\le2$: $f=x-2(x-1)=2-x$, so $f(1.5)=0.5$, $f(2)=0$. For $x\ge2$: $f=x-2(x-1)+(x-2)=0$. The result is a triangle of height 1 on $[0,2]$.
\end{example}
\begin{center}\begin{tikzpicture}
\begin{axis}[width=8cm,height=4.2cm,xmin=-1,xmax=3,ymin=-0.2,ymax=1.3,grid=major,grid style={GreyN!20},xlabel={$x$},ylabel={$f(x)$}]
\addplot[AccN,line width=1.3pt] coordinates{(-1,0)(0,0)(1,1)(2,0)(3,0)};
\end{axis}\end{tikzpicture}\\*{\small\color{GreyN}Figure 4: three hidden ReLU units make a bump.}\end{center}
Scaled and shifted bumps can be added to trace any continuous curve, like drawing a curve from many short straight segments. In two dimensions each hidden unit contributes one line, and the output combines them (``inside line 1 AND line 2 AND \dots''); six units can enclose a target inside a ring of clutter, and more units give a smoother region.

\subsection{The universal approximation theorem}
\begin{definition}{universal approximation (Cybenko 1989; Hornik 1991)}
A feed-forward network with one hidden layer of sufficiently many units with a suitable non-linear activation (sigmoid, ReLU and many others) can approximate any continuous function on a bounded region to any desired accuracy.
\end{definition}
So the answer to ``\emph{can} a network represent my problem?'' is almost always yes --- if it has a hidden layer with a non-linear activation. The fine print matters:
\begin{enumerate}
\item \textbf{How many units?} The theorem says enough units \emph{exist}; for complicated functions one hidden layer may need astronomically many.
\item \textbf{Can training find them?} It says nothing about whether gradient descent will find the right weights.
\item \textbf{Will it generalise?} Fitting the training data is not the same as predicting new data well (Lecture~1's overfitting and distribution shift).
\end{enumerate}

\subsection{Why depth: features of features}
A deep network computes a composition $f_3(f_2(f_1(x)))$. Each layer reuses everything below it, so a feature learned once (an edge, a ``low-altitude'' detector) serves many higher features. The number of distinct linear regions a ReLU network can carve out grows \emph{exponentially} with depth but only polynomially with width (Mont\'ufar et al., 2014): depth buys complexity cheaply. An analogy is a staff process --- sensor operators report contacts, analysts group them into formations, and the intelligence officer infers intent; each level builds on the one below.

Depth has a cost: deeper networks are harder to train, because gradients must pass back through many layers (Part~II, Section~\ref{sec:vanish}). ReLU activations, careful initialisation, normalisation layers and residual connections are the tools that made 100-layer networks practical.

\begin{pitfall}{five misconceptions about hidden layers}
\textbf{``More neurons in one layer solve everything.''} Without a hidden layer, each output is still one straight line.
\textbf{``More layers always help.''} Only if separated by non-linear activations, and only if the network can still be trained.
\textbf{``The activation is a detail.''} It is the entire source of the network's power beyond linear regression.
\textbf{``Hidden features are designed by engineers.''} They are learned by backpropagation from labels.
\textbf{``Universal approximation means any network will work.''} It guarantees existence, not size, trainability or generalisation.
\end{pitfall}

\section{Counting Parameters}
A dense layer from $n$ inputs to $m$ units has $nm$ weights and $m$ biases: $(n+1)m$ parameters. The XOR network has $(2+1)\times2=6$ hidden parameters and $2$ output weights (no output bias): 8 in all. The 2--2--1 network of Part~II has $(2+1)\times2+(2+1)\times1=9$ parameters --- the nine gradients computed there.
\begin{example}{a threat classifier}
A classifier with 20 track features, hidden layers of 64 and 32 ReLU units and a 5-class softmax output has $(20+1)\times64+(64+1)\times32+(32+1)\times5=1{,}344+2{,}080+165=3{,}589$ parameters. Lecture~7 shows how convolution keeps image networks small despite having thousands of inputs.
\end{example}

\section{Check Your Understanding (Part I)}
\begin{question}{1}
Give weights and a bias for a single step neuron that computes NAND (output 0 only when both inputs are 1).
\end{question}
\textbf{Answer.} NAND is NOT AND, so negate the AND neuron: $w=(-1,-1)$, $b=1.5$. Then $z=1.5,\,0.5,\,0.5,\,-0.5$ for $(0,0),(0,1),(1,0),(1,1)$, giving outputs $1,1,1,0$.

\begin{question}{2}
Which inequality in the XOR proof uses the $(0,0)$ row, and why is it needed?
\end{question}
\textbf{Answer.} Inequality (1), $b\le0$. Adding (2) and (3) only gives $w_1+w_2+b>-b$; we need $-b\ge0$, from (1), to conclude $w_1+w_2+b>0$ and contradict (4).

\begin{question}{3}
Collapse $W_1=\begin{pmatrix}2&0\\1&1\end{pmatrix}$, $b_1=0$, $W_2=(1\ \ 1)$, $b_2=-1$ into one layer.
\end{question}
\textbf{Answer.} $W=W_2W_1=(2+1,\ 0+1)=(3,\ 1)$ and $b=W_2b_1+b_2=-1$, so $\hat y=3x_1+x_2-1$.

\begin{question}{4}
In the XOR network, what does $\hat y$ become for $(0,0)$ if ReLU is replaced by the identity?
\end{question}
\textbf{Answer.} $z_1=0$, $z_2=-1$, so $\hat y=z_1-2z_2=0+2=2$ instead of 0. The linear network computes $-x_1-x_2+2$ and fails.

\begin{question}{5}
Name two advantages of learned feature maps over hand-engineered features.
\end{question}
\textbf{Answer.} Any two of: they are discovered automatically from labels; they are reusable across tasks (transfer learning); they compress many raw inputs into a few informative quantities; they can be visualised; they build hierarchically into more abstract features.

\begin{question}{6}
A single sigmoid neuron trained on XOR outputs 0.5 for every input. Is this a training failure that more epochs would fix?
\end{question}
\textbf{Answer.} No. Section~\ref{sec:logxor} shows the gradient is exactly zero at $w=0,b=0$, and the loss is convex, so this is the global optimum for a single neuron. Only a hidden layer can do better.

\part{Backpropagation by Hand}
\needspace{12\baselineskip}\section{The Network and Its Notation}
\needspace{8\baselineskip}\subsection{Architecture and given values}
The network has two inputs, one hidden layer of two sigmoid units, and one sigmoid output unit. It is trained with the binary cross-entropy (BCE) loss on a single example whose correct label is $y=1$.

\begin{center}
\begin{tikzpicture}[x=1cm,y=1cm,
 neu/.style={circle,draw=NavyN,line width=0.7pt,minimum size=1.3cm,inner sep=0pt,align=center,font=\small},
 inp/.style={neu,fill=GrnN!12,draw=GrnN},
 hid/.style={neu,fill=BluN!10,draw=BluN},
 outn/.style={neu,fill=BluN!25,draw=BluN},
 wl/.style={font=\footnotesize,fill=white,inner sep=1pt},
 bx/.style={draw=GreyN,fill=GreyN!8,rounded corners=2pt,font=\footnotesize,inner sep=2pt}]
\node[inp] (x1) at (0,1.5) {$x_1$\\1.0};
\node[inp] (x2) at (0,-1.5) {$x_2$\\0.5};
\node[hid] (h1) at (5,1.5) {$z^{(1)}_1|a^{(1)}_1$\\$\sigma$};
\node[hid] (h2) at (5,-1.5) {$z^{(1)}_2|a^{(1)}_2$\\$\sigma$};
\node[outn] (o) at (9.6,0) {$z^{(2)}|a^{(2)}$\\$\sigma$};
\node[draw=RedN,fill=RedN!8,rounded corners=2pt,align=center,font=\footnotesize] (L) at (12.4,0) {BCE loss\\$L(a^{(2)},y)$};
\node[bx] (b11) at (5,3.0) {$b^{(1)}_1=0$};
\node[bx] (b12) at (5,-3.0) {$b^{(1)}_2=0$};
\node[bx] (b2) at (9.6,1.7) {$b^{(2)}=0.1$};
\draw[arr] (x1)--node[wl,above=1pt]{$w^{(1)}_{11}=0.4$}(h1);
\draw[arr] (x2)--node[wl,below=1pt]{$w^{(1)}_{22}=0.6$}(h2);
\draw[arr] (x2)--node[wl,pos=0.72,sloped,above=1pt]{$w^{(1)}_{12}=-0.2$}(h1);
\draw[arr] (x1)--node[wl,pos=0.72,sloped,below=1pt]{$w^{(1)}_{21}=0.3$}(h2);
\draw[arr] (h1)--node[wl,above,sloped]{$w^{(2)}_1=0.7$}(o);
\draw[arr] (h2)--node[wl,below,sloped]{$w^{(2)}_2=-0.5$}(o);
\draw[arr] (o)--(L);
\draw[arr,dashed,draw=GreyN] (b11)--(h1);
\draw[arr,dashed,draw=GreyN] (b12)--(h2);
\draw[arr,dashed,draw=GreyN] (b2)--(o);
\node[font=\small\bfseries,text=GrnN] at (0,-3.9) {Input layer};
\node[font=\small\bfseries,text=BluN] at (5,-3.9) {Hidden layer (1)};
\node[font=\small\bfseries,text=BluN] at (9.6,-3.9) {Output layer (2)};
\node[font=\small\bfseries,text=RedN] at (12.4,-3.9) {Loss};
\end{tikzpicture}

\small Figure 5. The 2--2--1 network with every parameter shown on its connection. Dashed arrows are biases.
\end{center}

\[
x=\begin{pmatrix}x_1\\x_2\end{pmatrix}=\begin{pmatrix}1.0\\0.5\end{pmatrix},\quad y=1,\qquad
W^{(1)}=\begin{pmatrix}w^{(1)}_{11}&w^{(1)}_{12}\\[2pt]w^{(1)}_{21}&w^{(1)}_{22}\end{pmatrix}=\begin{pmatrix}0.4&-0.2\\0.3&0.6\end{pmatrix},\quad
b^{(1)}=\begin{pmatrix}0\\0\end{pmatrix},
\]
\[
w^{(2)}=\begin{pmatrix}w^{(2)}_1&w^{(2)}_2\end{pmatrix}=\begin{pmatrix}0.7&-0.5\end{pmatrix},\qquad b^{(2)}=0.1 .
\]

\needspace{8\baselineskip}\subsection{Notation: superscript for the layer, subscript for the index}
\begin{definition}{layered notation}
The \textbf{layer number} is always written as a parenthesised superscript and the \textbf{matrix or vector indices} as subscripts.
\begin{itemize}
\item $w^{(1)}_{jk}$ is the weight \emph{into} hidden unit $j$ \emph{from} input $x_k$: row $j$, column $k$ of $W^{(1)}$.
\item $b^{(1)}_j$, $z^{(1)}_j$ and $a^{(1)}_j$ are the bias, pre-activation and activation of hidden unit $j$.
\item $w^{(2)}_j$ is the weight from hidden unit $j$ to the output unit; $b^{(2)}$, $z^{(2)}$ and $a^{(2)}=\hat y$ belong to the output unit.
\item $\delta^{(l)}_j=\partial L/\partial z^{(l)}_j$ is the \textbf{error signal} at unit $j$ of layer $l$.
\end{itemize}
\end{definition}
The parentheses matter: $w^{(2)}$ means ``the weights of layer 2'', whereas $w^2$ would mean ``$w$ squared''. The order of the two subscripts also matters: the first names the receiving unit, the second the sending unit. With this convention one row of $W^{(1)}$ holds every weight entering one hidden unit, so the product $W^{(1)}x$ produces the vector of hidden pre-activations directly.

%=====================================================================
\needspace{12\baselineskip}\section{The Two Building Blocks}
\needspace{8\baselineskip}\subsection{The sigmoid activation}
\begin{definition}{logistic sigmoid}
$\displaystyle\sigma(z)=\frac{1}{1+e^{-z}}$ maps any real number $z$ into the open interval $(0,1)$.
\end{definition}

\begin{derivation}{the derivative of the sigmoid}
Write $\sigma(z)=(1+e^{-z})^{-1}$ and differentiate with the chain rule:
\[
\sigma'(z)=-(1+e^{-z})^{-2}\cdot(-e^{-z})=\frac{e^{-z}}{(1+e^{-z})^{2}}
=\underbrace{\frac{1}{1+e^{-z}}}_{\sigma(z)}\cdot\underbrace{\frac{e^{-z}}{1+e^{-z}}}_{1-\sigma(z)} .
\]
Hence $\sigma'(z)=\sigma(z)\bigl(1-\sigma(z)\bigr)$. The derivative can be computed from the activation alone --- no exponential needs to be re-evaluated in the backward pass.
\end{derivation}

\begin{center}
\begin{tikzpicture}
\begin{axis}[width=11cm,height=5.8cm,xmin=-6,xmax=6,ymin=0,ymax=1.35,grid=major,grid style={GreyN!20},xlabel={$z$},legend pos=north west,legend columns=2,legend style={font=\small}]
\addplot[BluN,line width=1.1pt,domain=-6:6,samples=150]{1/(1+exp(-x))};
\addplot[RedN,line width=1.1pt,domain=-6:6,samples=150]{exp(-x)/(1+exp(-x))^2};
\addplot[dashed,GreyN] coordinates{(-6,0.25)(6,0.25)};
\legend{$\sigma(z)$,$\sigma'(z)=\sigma(z)(1-\sigma(z))$}
\node[font=\footnotesize,text=RedN,anchor=south] at (axis cs:3.4,0.26) {maximum $\sigma'(0)=0.25$};
\end{axis}
\end{tikzpicture}

\small Figure 6. The sigmoid and its derivative. The slope never exceeds $0.25$.
\end{center}

\begin{keypoint}Because $\sigma'(z)\le 0.25$, every sigmoid layer that the error signal passes through multiplies it by at most one quarter. This is the origin of the \emph{vanishing-gradient problem} in deep sigmoid networks, and the main reason modern hidden layers use ReLU-type activations.\end{keypoint}

\needspace{8\baselineskip}\subsection{The binary cross-entropy loss}
\begin{definition}{binary cross-entropy}
For a label $y\in\{0,1\}$ and a predicted probability $a^{(2)}=\hat y\in(0,1)$,
\[
L=-\bigl[\,y\ln a^{(2)}+(1-y)\ln\bigl(1-a^{(2)}\bigr)\bigr].
\]
\end{definition}
The BCE is the negative log-likelihood of a Bernoulli label: if the network says ``the probability of class 1 is $\hat y$'', the probability it assigns to the label actually observed is $\hat y^{\,y}(1-\hat y)^{1-y}$, and $L$ is minus the logarithm of that. Minimising $L$ therefore maximises the likelihood of the training labels. When $y=1$ only the first term survives, $L=-\ln\hat y$; when $y=0$ only the second, $L=-\ln(1-\hat y)$.

\begin{center}
\begin{tikzpicture}
\begin{axis}[width=10cm,height=5.8cm,xmin=0,xmax=1,ymin=0,ymax=4,grid=major,grid style={GreyN!20},xlabel={prediction $\hat y=a^{(2)}$},ylabel={loss $L$},legend style={font=\small,at={(0.5,0.98)},anchor=north}]
\addplot[BluN,line width=1.1pt,domain=0.02:1,samples=120]{-ln(x)};
\addplot[GrnN,line width=1.1pt,domain=0:0.98,samples=120]{-ln(1-x)};
\addplot[only marks,mark=*,mark size=2pt,RedN] coordinates{(0.5447,0.6075)};
\draw[RedN] (axis cs:0.5447,0.6075)--(axis cs:0.45,2.2);
\node[font=\footnotesize,text=RedN,fill=white,inner sep=1pt,anchor=south] at (axis cs:0.45,2.2) {this example: $\hat y=0.5447$, $L=0.6075$};
\legend{$y=1$: $-\ln\hat y$,$y=0$: $-\ln(1-\hat y)$}
\end{axis}
\end{tikzpicture}

\small Figure 7. The BCE loss rises without limit as a prediction becomes confidently wrong.
\end{center}

\begin{defence}{the cost of confident errors}
Suppose a missile-warning classifier is shown a genuine threat ($y=1$). An undecided output $\hat y=0.5$ costs $L=-\ln 0.5=0.693$. A confident but wrong output $\hat y=0.01$ costs $L=-\ln 0.01=4.605$, more than six times as much. BCE therefore trains the network hardest on exactly the errors that are most dangerous operationally: the threats it is sure are harmless.
\end{defence}

%=====================================================================
\needspace{12\baselineskip}\section{The Forward Pass}
The forward pass moves left to right through Figure~5, computing each layer from the one before it. Every intermediate value is \emph{kept}, because the backward pass will need it.

\needspace{8\baselineskip}\subsection{F1 --- Hidden pre-activations}
Each hidden pre-activation is a weighted sum of the inputs plus a bias: $z^{(1)}=W^{(1)}x+b^{(1)}$, or unit by unit $z^{(1)}_j=\sum_k w^{(1)}_{jk}x_k+b^{(1)}_j$.
\begin{example}{F1}
\begin{align*}
z^{(1)}_1&=w^{(1)}_{11}x_1+w^{(1)}_{12}x_2+b^{(1)}_1=(0.4)(1.0)+(-0.2)(0.5)+0=0.4-0.1=\mathbf{0.300000}\\
z^{(1)}_2&=w^{(1)}_{21}x_1+w^{(1)}_{22}x_2+b^{(1)}_2=(0.3)(1.0)+(0.6)(0.5)+0=0.3+0.3=\mathbf{0.600000}
\end{align*}
In matrix form: $\begin{pmatrix}0.4&-0.2\\0.3&0.6\end{pmatrix}\begin{pmatrix}1.0\\0.5\end{pmatrix}+\begin{pmatrix}0\\0\end{pmatrix}=\begin{pmatrix}0.3\\0.6\end{pmatrix}$.
\end{example}

\needspace{8\baselineskip}\subsection{F2 --- Hidden activations}
\begin{example}{F2}
\begin{align*}
a^{(1)}_1&=\sigma(0.3)=\frac{1}{1+e^{-0.3}}=\frac{1}{1+0.740818}=\mathbf{0.574443}\\[2pt]
a^{(1)}_2&=\sigma(0.6)=\frac{1}{1+e^{-0.6}}=\frac{1}{1+0.548812}=\mathbf{0.645656}
\end{align*}
Quick check: near $z=0$, $\sigma(z)\approx0.5+z/4$, giving $0.575$ and $0.650$ --- close to the exact values, as expected for small $z$.
\end{example}

\needspace{8\baselineskip}\subsection{F3 --- Output pre-activation}
\begin{example}{F3}
\begin{align*}
z^{(2)}&=w^{(2)}_1a^{(1)}_1+w^{(2)}_2a^{(1)}_2+b^{(2)}=(0.7)(0.574443)+(-0.5)(0.645656)+0.1\\
&=0.402110-0.322828+0.1=\mathbf{0.179282}
\end{align*}
\end{example}
Notice the second term is negative: hidden unit~2 is active ($a^{(1)}_2=0.65$) but its outgoing weight is negative, so it \emph{pulls the output down}. This single observation explains a sign in the backward pass (Section~\ref{sec:b34}).

\needspace{8\baselineskip}\subsection{F4 --- Prediction}
\begin{example}{F4}
\[
a^{(2)}=\hat y=\sigma(0.179282)=\frac{1}{1+e^{-0.179282}}=\frac{1}{1+0.835871}=\mathbf{0.544701}
\]
\end{example}

\needspace{8\baselineskip}\subsection{F5 --- Loss}
\begin{example}{F5}
With $y=1$ the $(1-y)$ term is zero:
\[
L=-\bigl[\,1\cdot\ln(0.544701)+0\cdot\ln(0.455299)\,\bigr]=-\ln(0.544701)=\mathbf{0.607519}
\]
\end{example}

\begin{center}\renewcommand{\arraystretch}{1.25}
\begin{tabular}{lrlr}
\toprule
\multicolumn{4}{c}{\textbf{Table 1. Forward-pass results}}\\
\midrule
$z^{(1)}_1$&0.300000&$a^{(1)}_1$&0.574443\\
$z^{(1)}_2$&0.600000&$a^{(1)}_2$&0.645656\\
$z^{(2)}$&0.179282&$a^{(2)}=\hat y$&0.544701\\
\midrule
\multicolumn{4}{c}{Loss $L=0.607519$}\\
\bottomrule
\end{tabular}
\end{center}
The network gives a 54.5\% probability to an example whose true label is 1: barely better than a coin-flip. The rest of this text finds out how each parameter should change to raise that probability.

%=====================================================================
\needspace{12\baselineskip}\section{The Backward Pass}
\needspace{8\baselineskip}\subsection{The chain rule and why ``backpropagation''}
The loss depends on a weight in the first layer only through a chain of intermediate quantities. For $w^{(1)}_{11}$ the chain is
\[
w^{(1)}_{11}\ \xrightarrow{\ \times x_1\ }\ z^{(1)}_1\ \xrightarrow{\ \sigma\ }\ a^{(1)}_1\ \xrightarrow{\ \times w^{(2)}_1\ }\ z^{(2)}\ \xrightarrow{\ \sigma\ }\ a^{(2)}\ \xrightarrow{\ \text{BCE}\ }\ L ,
\]
and the chain rule multiplies the local derivative of every link:
\[
\frac{\partial L}{\partial w^{(1)}_{11}}=
\underbrace{\frac{\partial L}{\partial a^{(2)}}\frac{\partial a^{(2)}}{\partial z^{(2)}}}_{\delta^{(2)}}\;
\frac{\partial z^{(2)}}{\partial a^{(1)}_1}\;
\frac{\partial a^{(1)}_1}{\partial z^{(1)}_1}\;
\frac{\partial z^{(1)}_1}{\partial w^{(1)}_{11}}
=\underbrace{\delta^{(2)}\,w^{(2)}_1\,\sigma'(z^{(1)}_1)}_{\delta^{(1)}_1}\;x_1 .
\]
\begin{keypoint}Backpropagation is the chain rule \emph{plus reuse}. The first two factors, $\delta^{(2)}$, are common to every parameter in the network, so they are computed once. The first four factors, $\delta^{(1)}_1$, are common to every weight entering hidden unit~1, so they too are computed once. Working from the loss backwards guarantees that each shared factor is available when it is needed. This is why the cost of computing \emph{all} gradients is only about twice the cost of one forward pass.\end{keypoint}

\needspace{8\baselineskip}\subsection{B1 --- The output error signal}
\begin{derivation}{$\delta^{(2)}=a^{(2)}-y$ for a sigmoid output with BCE loss}
Differentiate the loss with respect to the prediction:
\[
\frac{\partial L}{\partial a^{(2)}}=-\frac{y}{a^{(2)}}+\frac{1-y}{1-a^{(2)}}
=\frac{-y(1-a^{(2)})+(1-y)a^{(2)}}{a^{(2)}(1-a^{(2)})}
=\frac{a^{(2)}-y}{a^{(2)}\bigl(1-a^{(2)}\bigr)} .
\]
The sigmoid contributes $\partial a^{(2)}/\partial z^{(2)}=a^{(2)}(1-a^{(2)})$. Multiplying,
\[
\delta^{(2)}=\frac{\partial L}{\partial z^{(2)}}=\frac{a^{(2)}-y}{a^{(2)}(1-a^{(2)})}\cdot a^{(2)}(1-a^{(2)})=a^{(2)}-y .
\]
\end{derivation}
\begin{example}{B1}
\[
\delta^{(2)}=a^{(2)}-y=0.544701-1=\mathbf{-0.455299}
\]
\end{example}
\begin{pitfall}{using squared error with a sigmoid output}
Had we used $L=\tfrac12(\hat y-y)^2$, the sigmoid factor would \emph{not} cancel: $\delta^{(2)}=(\hat y-y)\,\hat y(1-\hat y)=(-0.455299)(0.544701)(0.455299)=-0.112915$, four times weaker. Worse, if the network were confidently wrong ($\hat y\approx0$ for $y=1$), the factor $\hat y(1-\hat y)\approx0$ would almost silence the gradient precisely when learning is most needed. Sigmoid outputs should be paired with BCE.
\end{pitfall}

\needspace{8\baselineskip}\subsection{B2 --- Output-layer gradients}
Because $z^{(2)}=w^{(2)}_1a^{(1)}_1+w^{(2)}_2a^{(1)}_2+b^{(2)}$, we have $\partial z^{(2)}/\partial w^{(2)}_j=a^{(1)}_j$ and $\partial z^{(2)}/\partial b^{(2)}=1$. Hence $\partial L/\partial w^{(2)}_j=\delta^{(2)}a^{(1)}_j$ and $\partial L/\partial b^{(2)}=\delta^{(2)}$.
\begin{example}{B2}
\begin{align*}
\frac{\partial L}{\partial w^{(2)}_1}&=\delta^{(2)}a^{(1)}_1=(-0.455299)(0.574443)=\mathbf{-0.261543}\\[2pt]
\frac{\partial L}{\partial w^{(2)}_2}&=\delta^{(2)}a^{(1)}_2=(-0.455299)(0.645656)=\mathbf{-0.293967}\\[2pt]
\frac{\partial L}{\partial b^{(2)}}&=\delta^{(2)}=\mathbf{-0.455299}
\end{align*}
\end{example}
All three are negative. Gradient \emph{descent} subtracts the gradient, so it will \emph{increase} $w^{(2)}_1$, $w^{(2)}_2$ and $b^{(2)}$; each of these increases raises $z^{(2)}$ and therefore pushes $\hat y$ towards the target 1. The rule to remember is \textbf{gradient = (error signal at the receiving unit) $\times$ (activation on the sending side)}; a bias always sends the value 1.

\needspace{8\baselineskip}\subsection{B3 and B4 --- Error signals at the hidden layer}\label{sec:b34}
The error reaches hidden unit $j$ through the single weight $w^{(2)}_j$ and is then scaled by the slope of that unit's sigmoid:
\[
\delta^{(1)}_j=\frac{\partial L}{\partial z^{(1)}_j}=\delta^{(2)}\,w^{(2)}_j\,\sigma'(z^{(1)}_j),\qquad \sigma'(z^{(1)}_j)=a^{(1)}_j\bigl(1-a^{(1)}_j\bigr).
\]
\begin{example}{B3 --- sigmoid slopes}
\begin{align*}
\sigma'(z^{(1)}_1)&=(0.574443)(1-0.574443)=(0.574443)(0.425557)=\mathbf{0.244458}\\
\sigma'(z^{(1)}_2)&=(0.645656)(1-0.645656)=(0.645656)(0.354344)=\mathbf{0.228784}
\end{align*}
\end{example}
\begin{example}{B4 --- hidden error signals}
\begin{align*}
\delta^{(1)}_1&=(-0.455299)(0.7)(0.244458)=(-0.318709)(0.244458)=\mathbf{-0.077911}\\
\delta^{(1)}_2&=(-0.455299)(-0.5)(0.228784)=(0.227650)(0.228784)=\mathbf{0.052083}
\end{align*}
\end{example}
The positive sign of $\delta^{(1)}_2$ is the backward-pass echo of the observation made in F3. Hidden unit~2 feeds the output through a \emph{negative} weight, so increasing its activity lowers $\hat y$ --- the wrong direction for a positive example. The gradient therefore tells unit~2 to become less active. Note also the size of the signals: $|\delta^{(2)}|=0.455$ but $|\delta^{(1)}_j|\le0.078$. One weight smaller than 1 and one sigmoid slope at most 0.25 have reduced the signal about six-fold in a single layer.

\needspace{8\baselineskip}\subsection{B5 --- Hidden-layer gradients}
Since $z^{(1)}_j=\sum_k w^{(1)}_{jk}x_k+b^{(1)}_j$, we have $\partial z^{(1)}_j/\partial w^{(1)}_{jk}=x_k$ and $\partial z^{(1)}_j/\partial b^{(1)}_j=1$:
\[
\frac{\partial L}{\partial w^{(1)}_{jk}}=\delta^{(1)}_j\,x_k,\qquad \frac{\partial L}{\partial b^{(1)}_j}=\delta^{(1)}_j .
\]
\begin{example}{B5}
\begin{align*}
\frac{\partial L}{\partial w^{(1)}_{11}}&=(-0.077911)(1.0)=\mathbf{-0.077911}, &
\frac{\partial L}{\partial w^{(1)}_{12}}&=(-0.077911)(0.5)=\mathbf{-0.038956},\\
\frac{\partial L}{\partial w^{(1)}_{21}}&=(0.052083)(1.0)=\mathbf{0.052083}, &
\frac{\partial L}{\partial w^{(1)}_{22}}&=(0.052083)(0.5)=\mathbf{0.026041},\\
\frac{\partial L}{\partial b^{(1)}_{1}}&=\mathbf{-0.077911}, &
\frac{\partial L}{\partial b^{(1)}_{2}}&=\mathbf{0.052083}.
\end{align*}
In matrix form the whole weight gradient is one outer product:
\[
\nabla_{W^{(1)}}L=\delta^{(1)}x^{\mathsf T}=\begin{pmatrix}-0.077911\\0.052083\end{pmatrix}\begin{pmatrix}1.0&0.5\end{pmatrix}
=\begin{pmatrix}-0.077911&-0.038956\\0.052083&0.026041\end{pmatrix}.
\]
\end{example}

\needspace{8\baselineskip}\subsection{The backward flow at a glance}
\begin{center}
\begin{tikzpicture}[x=1cm,y=1cm,
 neu/.style={circle,draw=RedN,line width=0.7pt,minimum size=1.3cm,inner sep=0pt,align=center,font=\small},
 bar/.style={{Stealth[length=4pt]}-,thick,draw=RedN},
 wl/.style={font=\footnotesize,fill=white,inner sep=1pt,text=RedN}]
\node[neu,draw=GrnN,fill=GrnN!12] (x1) at (0,1.5) {$x_1$\\1.0};
\node[neu,draw=GrnN,fill=GrnN!12] (x2) at (0,-1.5) {$x_2$\\0.5};
\node[neu,fill=RedN!8] (h1) at (5,1.5) {$\delta^{(1)}_1$\\$-0.0779$};
\node[neu,fill=RedN!8] (h2) at (5,-1.5) {$\delta^{(1)}_2$\\$0.0521$};
\node[neu,fill=RedN!20] (o) at (9.8,0) {$\delta^{(2)}$\\$-0.4553$};
\draw[bar] (x1)--node[wl,above=1pt]{$-0.0779$}(h1);
\draw[bar] (x2)--node[wl,below=1pt]{$0.0260$}(h2);
\draw[bar] (x2)--node[wl,pos=0.72,sloped,above=1pt]{$-0.0390$}(h1);
\draw[bar] (x1)--node[wl,pos=0.72,sloped,below=1pt]{$0.0521$}(h2);
\draw[bar] (h1)--node[wl,above,sloped]{$-0.2615$}(o);
\draw[bar] (h2)--node[wl,below,sloped]{$-0.2940$}(o);
\end{tikzpicture}

\small Figure 8. Backward flow. Nodes carry $\delta=\partial L/\partial z$; edges carry $\partial L/\partial w$ for that connection.
\end{center}

\begin{center}\renewcommand{\arraystretch}{1.25}
\begin{tabular}{lr@{\qquad\qquad}lr}
\toprule
\multicolumn{4}{c}{\textbf{Table 2. All gradients}}\\
\midrule
\textbf{Parameter}&$\partial L/\partial(\cdot)$&\textbf{Parameter}&$\partial L/\partial(\cdot)$\\
\midrule
$w^{(1)}_{11}$&$-0.077911$&$w^{(2)}_1$&$-0.261543$\\
$w^{(1)}_{12}$&$-0.038956$&$w^{(2)}_2$&$-0.293967$\\
$w^{(1)}_{21}$&$0.052083$&$b^{(2)}$&$-0.455299$\\
$w^{(1)}_{22}$&$0.026041$&$\delta^{(2)}$&$-0.455299$\\
$b^{(1)}_1$&$-0.077911$&$\delta^{(1)}_1$&$-0.077911$\\
$b^{(1)}_2$&$0.052083$&$\delta^{(1)}_2$&$0.052083$\\
\bottomrule
\end{tabular}
\end{center}

%=====================================================================
\needspace{12\baselineskip}\section{Verifying the Gradients}
\needspace{8\baselineskip}\subsection{Numerical gradient check}
A hand-derived or hand-coded gradient should always be checked against a finite-difference estimate,
\[
\frac{\partial L}{\partial\theta}\approx\frac{L(\theta+\varepsilon)-L(\theta-\varepsilon)}{2\varepsilon}.
\]
\begin{example}{central-difference check of $\partial L/\partial w^{(1)}_{11}$ with $\varepsilon=0.01$}
Re-running the forward pass with $w^{(1)}_{11}=0.41$ gives $L=0.6067406$; with $w^{(1)}_{11}=0.39$ it gives $L=0.6082988$. Then
\[
\frac{0.6067406-0.6082988}{0.02}=\frac{-0.0015582}{0.02}=-0.077911,
\]
in agreement with the backpropagated value to six decimal places. The same check for all nine parameters reproduces every entry of Table~2.
\end{example}
Deep-learning libraries compute gradients automatically, but gradient checking remains the standard test when a new layer or loss is implemented by hand.

\needspace{8\baselineskip}\subsection{One gradient-descent step}
The problem does not specify a learning rate; for illustration take $\eta=0.5$ and apply $\theta\leftarrow\theta-\eta\,\partial L/\partial\theta$ to every parameter \emph{simultaneously}, using gradients computed from the old values.
\begin{example}{update with $\eta=0.5$}
\begin{align*}
W^{(1)}&\leftarrow\begin{pmatrix}0.4-0.5(-0.077911)&-0.2-0.5(-0.038956)\\0.3-0.5(0.052083)&0.6-0.5(0.026041)\end{pmatrix}=\begin{pmatrix}0.438956&-0.180522\\0.273959&0.586979\end{pmatrix}\\[3pt]
b^{(1)}&\leftarrow(0.038956,\ -0.026041),\qquad w^{(2)}\leftarrow(0.830772,\ -0.353017),\qquad b^{(2)}\leftarrow0.327650
\end{align*}
Repeating the forward pass with these values gives $\hat y=0.645519$ and $L=0.437701$. The prediction has moved towards the target and the loss has fallen from $0.6075$ to $0.4377$, confirming that the gradients point in the right direction.
\end{example}

\begin{center}
\begin{tikzpicture}
\begin{axis}[width=11cm,height=5.8cm,xmin=0,xmax=30,ymin=0,ymax=0.65,grid=major,grid style={GreyN!20},xlabel={gradient-descent step},ylabel={loss $L$}]
\addplot[BluN,line width=1.1pt,mark=*,mark size=1pt] coordinates{(0,0.6075) (1,0.4377) (2,0.3326) (3,0.2638) (4,0.2164) (5,0.1822) (6,0.1567) (7,0.1369) (8,0.1213) (9,0.1087) (10,0.0983) (11,0.0896) (12,0.0822) (13,0.0759) (14,0.0705) (15,0.0657) (16,0.0615) (17,0.0578) (18,0.0545) (19,0.0515) (20,0.0488) (21,0.0464) (22,0.0442) (23,0.0422) (24,0.0403) (25,0.0386) (26,0.0371) (27,0.0356) (28,0.0343) (29,0.0330) (30,0.0318)};
\addplot[only marks,mark=*,mark size=2pt,RedN] coordinates{(1,0.4377)};
\node[font=\footnotesize,text=RedN,anchor=west] at (axis cs:2,0.45) {after the step worked above};
\end{axis}
\end{tikzpicture}

\small Figure 9. Repeating the same step on the same example ($\eta=0.5$): after 10 steps $\hat y=0.906$, after 30 steps $\hat y=0.969$.
\end{center}

%=====================================================================

\needspace{8\baselineskip}\subsection{Why gradients vanish in deep sigmoid networks}\label{sec:vanish}
The hidden error signal $\delta^{(1)}_j=\delta^{(2)}w^{(2)}_j\sigma'(z^{(1)}_j)$ shows the general pattern: each layer the error passes back through multiplies it by a weight and by the derivative of that layer's activation. With $L$ sigmoid layers the signal reaching the first layer contains a product of $L$ factors $\sigma'(z)\le0.25$. Even with weights of size 1, ten layers shrink it by at least $0.25^{10}\approx9.5\times10^{-7}$: the early layers barely learn. This is the \textbf{vanishing-gradient problem}, and it is why the ReLU of Part~I (derivative exactly 1 for active units) replaced the sigmoid in hidden layers, and why very deep networks add \emph{residual connections} $y=x+F(x)$, whose derivative contains a direct term 1 that carries the gradient past each block. Lecture~9 meets residual connections again in every transformer block.

\needspace{12\baselineskip}\section{Defence and Wargaming Perspective}
\begin{defence}{a neural threat classifier}
Replace the two inputs by two radar-track features --- say, normalised closing speed and a normalised IFF-response score --- and read the output as $P(\text{hostile}\mid x)$. The weights then have an operational meaning. A negative weight such as $w^{(1)}_{12}=-0.2$ means the second feature \emph{inhibits} the first hidden detector, as a strong IFF reply should inhibit a ``threat'' detector. The gradient signs found above tell us how one hostile training track reshapes the classifier: the detector that supports a hostile call (unit~1) is strengthened, while the detector that argues against it (unit~2) is weakened.
\end{defence}
\begin{defence}{why vanishing gradients matter operationally}
Deep sigmoid networks trained on sensor data can stall: the early layers, which learn the most basic features of the signal, receive gradients shrunk by a factor of 4 or more per layer and barely change. A model that appears to train (its output layer adapts) may in fact be relying on poorly-learned features that do not transfer to new theatres --- a direct link to the distribution-shift failures discussed in Lecture~1.
\end{defence}

%=====================================================================
\needspace{12\baselineskip}\section{Common Mistakes}
\begin{pitfall}{forgetting $\sigma'$ at a hidden unit}
$\delta^{(1)}_j=\delta^{(2)}w^{(2)}_j$ is wrong. The error must also pass through the hidden sigmoid: multiply by $a^{(1)}_j(1-a^{(1)}_j)$.
\end{pitfall}
\begin{pitfall}{adding $\sigma'$ at the output when using BCE}
With a sigmoid output and BCE loss the factor $a^{(2)}(1-a^{(2)})$ has already cancelled. Writing $\delta^{(2)}=(\hat y-y)\sigma'(z^{(2)})$ counts it twice.
\end{pitfall}
\begin{pitfall}{swapping the indices}
$w^{(1)}_{jk}$ goes \emph{to} unit $j$ \emph{from} input $k$, and its gradient is $\delta^{(1)}_j x_k$. Writing $\delta^{(1)}_k x_j$ transposes the gradient matrix.
\end{pitfall}
\begin{pitfall}{updating weights in the middle of the backward pass}
$\delta^{(1)}_j$ uses the \emph{old} value of $w^{(2)}_j$. Compute every gradient first, then update every parameter together.
\end{pitfall}
\begin{pitfall}{the sign of the step}
Gradient descent is $\theta-\eta\nabla L$. Adding the gradient climbs the loss surface. After any hand calculation, re-run the forward pass: for a small enough $\eta$ the loss must fall.
\end{pitfall}

%=====================================================================
\needspace{12\baselineskip}\section{Summary of Equations}
\begin{center}\small\renewcommand{\arraystretch}{1.6}
\begin{tabular}{ll}
\toprule
\textbf{Forward} & $z^{(1)}=W^{(1)}x+b^{(1)},\quad a^{(1)}=\sigma(z^{(1)}),\quad z^{(2)}=w^{(2)}a^{(1)}+b^{(2)},\quad a^{(2)}=\sigma(z^{(2)})$\\
\textbf{Loss} & $L=-\bigl[y\ln a^{(2)}+(1-y)\ln(1-a^{(2)})\bigr]$\\
\textbf{Output error} & $\delta^{(2)}=a^{(2)}-y$\\
\textbf{Hidden error} & $\delta^{(1)}_j=\delta^{(2)}\,w^{(2)}_j\,a^{(1)}_j(1-a^{(1)}_j)$\\
\textbf{Gradients} & $\partial L/\partial w^{(2)}_j=\delta^{(2)}a^{(1)}_j,\quad \partial L/\partial b^{(2)}=\delta^{(2)},\quad \partial L/\partial w^{(1)}_{jk}=\delta^{(1)}_jx_k,\quad \partial L/\partial b^{(1)}_j=\delta^{(1)}_j$\\
\textbf{Update} & $\theta\leftarrow\theta-\eta\,\partial L/\partial\theta$\\
\bottomrule
\end{tabular}
\end{center}

%=====================================================================
\needspace{12\baselineskip}\section{Check Your Understanding (Part II)}
\begin{question}{1}
If the label were $y=0$ instead of $1$, with everything else unchanged, what would $\delta^{(2)}$ be, and in which direction would $b^{(2)}$ move?
\end{question}
\textbf{Answer.} $\delta^{(2)}=a^{(2)}-y=0.544701-0=0.544701$. The gradient $\partial L/\partial b^{(2)}=0.544701$ is positive, so gradient descent \emph{decreases} $b^{(2)}$, lowering $\hat y$ towards 0.

\begin{question}{2}
If $w^{(2)}_2$ were $+0.5$ instead of $-0.5$, what would happen to the weights entering hidden unit~2?
\end{question}
\textbf{Answer.} $\delta^{(1)}_2=(-0.455299)(+0.5)(0.228784)=-0.052083$, now negative, so $\partial L/\partial w^{(1)}_{21}$ and $\partial L/\partial w^{(1)}_{22}$ become negative and gradient descent \emph{increases} both weights. Unit~2 would now support the positive output, so the network makes it more active --- the same treatment unit~1 receives.

\begin{question}{3}
Why is $|\partial L/\partial w^{(1)}_{12}|$ exactly half of $|\partial L/\partial w^{(1)}_{11}|$?
\end{question}
\textbf{Answer.} Both weights enter the same hidden unit and so share the factor $\delta^{(1)}_1$; their gradients are $\delta^{(1)}_1x_2$ and $\delta^{(1)}_1x_1$, and $x_2=0.5$ is half of $x_1=1.0$. Weights entering one unit always have gradients in proportion to their inputs.

\begin{question}{4}
What is the largest possible value of $\sigma'(z)$, and why does it matter for deep networks?
\end{question}
\textbf{Answer.} $\sigma'(z)=\sigma(z)(1-\sigma(z))$ is largest when $\sigma(z)=0.5$, i.e.\ $z=0$, where it equals $0.25$. Each sigmoid layer therefore scales the backward signal by at most $0.25$ (times the weight); across ten layers the factor is at most $0.25^{10}\approx10^{-6}$, so early layers learn extremely slowly --- the vanishing-gradient problem.

\end{document}
